\subsection{Do-Calculus}
\totalscore{9}

\textbf{TODO: this exercise is wrong! The answer is wrong, and it could be that the correct answer is that the target
quantity is not even identifiable from the given quantities (but we are not sure about this).}

\begin{figure}[ht]
\centering
\begin{tikzpicture}[scale=1, transform shape]
    \node[ndout] (X) at (0,0) {$X$};
    \node[ndout] (Y) at (6,0) {$Y$};
    \node[ndout] (Z) at (3,0) {$Z$};
    \node[ndout] (R) at (3,2) {$R$};
    \node[ndout] (L) at (0,2) {$L$};
    \node[ndout] (S) at (3,-2) {$S$};
    \node[ndout] (W) at (6,2) {$W$};
    \draw[arout] (Z) to (S);
    \draw[arout] (W) to (Y);
    \draw[arout] (X) to (Z);
    \draw[arout] (Z) to (Y);
    \draw[arout] (R) to (Y);
    \draw[arout] (R) to (X);
    \draw[arout] (L) to (X);
    \draw[arlat, bend left] (L) to (R);
    \draw[arlat, bend left] (R) to (W);
\end{tikzpicture}
\caption{Acyclic Directed Mixed Graph (ADMG).}
\label{fig:cadmg-1}
\end{figure}

\noindent Consider the observational graph of a causal Bayesian network with latent variables in Figure \ref{fig:cadmg-1}.
%We assume that the variables $X,Y,Z,W,L,S,R$ are observable with probability distribution $P(X,Y,Z,W,L,S,R)$.
To avoid measure-theoretic subtleties, we assume that all occuring Markov kernels are discrete and have strictly positive mass functions.

In this exercise we are interested in estimating the causal effect of $X$ and $W$ on $Y$, that is, the interventional distribution $P(Y|\doit(X,W))$. However, we are only able to perform certain interventional experiments:
\begin{enumerate}[label=(\arabic*)]
    \item We can perform experiments where we intervene on $W$, but this comes, unfortunately, with certain selection bias, $S=s$. In other words, we are only able to produce samples from the conditional interventional distribution: $P(X,Y,Z,L,R|\doit(W),S=s)$ for a fixed value $s$.
    \item We can perform additional unbiased experiments where we intervene on $W$, but there we can only measure the output of $X$ and $Z$. This means that we can also produce samples from the (unbiased, unconditional)  interventional distribution $P(X,Z|\doit(W))$.
\end{enumerate}

Use the 3 rules of do-calculus together with the rules of probability (e.g.\ chain rule) to derive an expression for the Markov kernel $P(Y|\doit(X,W))$ in terms of $P(X,Y,Z,L,R|S=s,\doit(W))$ and $P(X,Z|\doit(W))$. This expression can contain combinations of marginals, conditionals, products and compositions of the latter two Markov kernels. 
Provide all details of your arguments.
When you use d-separation statements clearly indicate in which graph this should hold and draw the graph if it deviates from the one depicted in Figure \ref{fig:cadmg-1}.
%The claimed d-separation statements itself need to be proven as well, by clearly indicating which paths need to be considered and why and where they are d-blocked.
The claimed d-separation statements do not need to be proven (but should be correct, of course).

\emph{Hint: There exists a solution, where the 3 rules of the do-calculus rules are each applied once. Maybe focus on the chain rule of the form $P(Y|\doit(X,W)) = P(Y|U,\doit(X,W)) \circ P(U|\doit(X,W))$ for suitable variables $U$ first and then remove/introduce expressions using the do-calculus rules.}


\answer{%
Since all occuring Markov kernels contain $\doit(W)$ we can restrict to the following CADMG:
\points{1pt for drawing the correct interventional graph with intervention node in the following.}

\centerline{
\begin{tikzpicture}[scale=1, transform shape]
    \node[ndout] (X) at (0,0) {$X$};
    \node[ndout] (Y) at (6,0) {$Y$};
    \node[ndout] (Z) at (3,0) {$Z$};
    \node[ndout] (R) at (3,2) {$R$};
    \node[ndout] (L) at (0,2) {$L$};
    \node[ndout] (S) at (3,-2) {$S$};
    \node[ndint] (W) at (6,2) {$W$};
    \draw[arout] (Z) to (S);
    \draw[arout] (W) to (Y);
    \draw[arout] (X) to (Z);
    \draw[arout] (Z) to (Y);
    \draw[arout] (R) to (Y);
    \draw[arout] (R) to (X);
    \draw[arout] (L) to (X);
    \draw[arlat, bend left] (L) to (R);
\end{tikzpicture}
}

We can then add the intervention node $I_X$ to the graph and also marginalize out $L$ and $R$. The latter is not important for the results but makes the graph easier to parse.

\centerline{
\begin{tikzpicture}[scale=1, transform shape]
    \node[ndout] (X) at (0,0) {$X$};
    \node[ndint] (IX) at (-2,1) {$I_X$};
    \node[ndout] (Y) at (6,0) {$Y$};
    \node[ndout] (Z) at (3,0) {$Z$};
    \node[ndout] (S) at (3,-2) {$S$};
    \node[ndint] (W) at (6,2) {$W$};
    \draw[arout] (Z) to (S);
    \draw[arout] (W) to (Y);
    \draw[arout] (X) to (Z);
    \draw[arout] (Z) to (Y);
    \draw[arint] (IX) to (X);
    \draw[arlat, bend left] (X) to (Y);
\end{tikzpicture}
}



\begin{enumerate}
    \item By the chain rule we get:
        \begin{align*}
            P(Y|\doit(X,W)) &= P(Y|Z,\doit(X,W)) \circ P(Z|\doit(X,W)).
        \end{align*}
        \points{1pt for the chain rule.}
          \item To apply the 2nd rule of do-calculus to $P(Z|\doit(X,W))$ we check the d-separation in the depicted graph $G_{\doit(I_X,W)}$:
         \begin{align*}
             Z \Perp^d_{G_{\doit(I_X,W)}} I_X \given X,W.
         \end{align*}
         The 2nd rule of do-calculus then shows the equation:
         \begin{align*}
             P(Z|\doit(X,W)) & = P(Z|X,\doit(W)).
         \end{align*}
         \points{1pt for correct d-separation statement. 1pt for correct application of do-calculus rule.
         }
     \item To apply the 3rd rule of do-calculus to $P(Y|Z,\doit(X,W))$ we check the d-separation in the depicted graph $G_{\doit(I_X,W)}$: 
         \begin{align*}
             Y \Perp^d_{G_{\doit(I_X,W)}} I_X \given Z,W.
         \end{align*}
         \textbf{TODO: this is wrong, because conditioning on $Z$ opens up its ancestor collider $X$.}
         The 3rd rule of do-calculus then shows the equation:
         \begin{align*}
             P(Y|Z,\doit(X,W)) & = P(Y|Z,\doit(W)).
         \end{align*}
         \points{1pt for correct d-separation statement.
         1pt for correct application of do-calculus rule.
         }
     \item To apply the 1st rule of do-calculus to $P(Y|Z,\doit(W))$ we check the d-separation in the depicted graph $G_{\doit(I_X,W)}$ (or $G_{\doit(W)}$) : 
         \begin{align*}
             Y \Perp^d_{G_{\doit(I_X,W)}} S \given Z,W.
         \end{align*}
         The 1st rule of do-calculus then shows the equation:
         \begin{align*}
             P(Y|Z,\doit(W)) & = P(Y|Z,S,\doit(W)).
         \end{align*}
         Since the left hand side is not dependent on $S$ we can evaluate at $S=s$ and get:
         \begin{align*}
             P(Y|Z,\doit(W)) & = P(Y|Z,S=s,\doit(W)).
         \end{align*}
         \points{1pt for correct d-separation statement. 
           1pt for correct application of do-calculus rule.  
         }
     \item Putting all points together we get:
        \begin{align*}
            P(Y|\doit(X,W)) &= P(Y|Z,\doit(X,W)) \circ P(Z|\doit(X,W)) \\
                            &= P(Y|Z,S=s,\doit(W)) \circ P(Z|X,\doit(W)).
        \end{align*}
        \points{1pt for putting it all together.}
\end{enumerate}
}
